\section{Related Work}\label{sec:related}
\subsection{Reduced-basis surrogates with learned coefficients}
Non-intrusive reduced-order models regress the weights of POD modes on the parameters. Hesthaven and Ubbiali used neural networks for this~\cite{hesthaven2018}, Bhattacharya et al.\ analysed PCA-based neural operators~\cite{bhattacharya2021}, and Guo and Hesthaven used Gaussian-process regression for nonlinear structural analysis~\cite{guo2018}. B\={e}rzi\c{n}\v{s} et al.\ compared radial-basis-function, Gaussian-process and neural-network regression inside one POD framework for extruder flows and reported several advantages of the Gaussian process~\cite{berzins2020}; we read only the abstract, and it is the closest precedent for our finding on smooth problems. Mora et al.\ proposed a hybrid Gaussian-process/neural operator framework~\cite{mora2025}.

\subsection{Neural networks, operators and nonlinear manifolds}
DeepONet~\cite{lu2021}, the Fourier neural operator~\cite{li2021} and the neural-operator framework~\cite{kovachki2023} learn maps between function spaces; autoencoders learn nonlinear latent manifolds for reduced models~\cite{lee2020,fresca2022,franco2023}. Our parameters are finite-dimensional, so we study direct and bottleneck networks rather than operator-input architectures. Whether a small linear basis can work depends on the decay of the Kolmogorov $n$-width~\cite{cohen2015}, which is slow for wave problems with discontinuous data~\cite{greif2019}.

\subsection{Kernel and Gaussian-process methods}
Gaussian processes are the standard low-dimensional regressor~\cite{rasmussen2006}. Batlle et al.\ develop a kernel framework with a priori error analysis and report it competitive with DeepONet and the Fourier neural operator~\cite{batlle2024}. A position paper argues that many neural-surrogate successes sit on low-dimensional manifolds where any competent reduced model interpolates well~\cite{duraisamy2026}.

\subsection{Fair comparison, tuning and reporting}
McGreivy and Hakim found that 79\% of 76 reviewed articles compared with weak baselines~\cite{mcgreivy2024}. De Hoop et al.\ study the cost--accuracy trade-off of neural architectures~\cite{dehoop2022}; Westermann et al.\ compare neural and polynomial surrogates by Pareto frontiers and find no universal winner~\cite{westermann2026}; a Common Task Framework has been proposed to evaluate scientific machine-learning methods against curated baselines~\cite{wyder2025}. Random search is an efficient alternative to grid search~\cite{bergstra2012}; adaptive early-stopping schemes such as Hyperband~\cite{li2018hyperband} and software such as Optuna~\cite{akiba2019} reduce tuning cost; reporting results as a function of tuning budget and accounting for run-to-run variance have been advocated~\cite{dodge2019,bouthillier2021}.

\subsection{Surrogate selection, learning curves and software}
Williams and Cremaschi compared eight surrogate forms on generated data to look for rules for choosing a form~\cite{williams2019}. Learning curves, the error as a function of training-set size, have been surveyed for decision making in supervised learning~\cite{mohr2024}, and error scaling laws have been fitted across data scales in deep learning~\cite{rosenfeld2020}. General-purpose libraries for model-order reduction~\cite{milk2016} and uncertainty quantification~\cite{marelli2014} exist; we did not compare their features with ours, and the finite-element library \texttt{scikit-fem}~\cite{gustafsson2020} is used here only to generate independent test problems. PDEBench~\cite{takamoto2022} provides large simulated benchmarks for scientific machine learning; we did not use it (\cref{sec:discussion}).

\subsection{Comparison with prior work}
\Cref{tab:related} lists what each closest work contains, judged from the parts we read (mostly abstracts; Westermann et al.\ in large part).
\begin{table}[t]\centering\caption{Closest related work. A dash means that we did not find the property in the parts we read (mostly abstracts); it does not prove absence.}\label{tab:related}\scriptsize\setlength{\tabcolsep}{2pt}\resizebox{\textwidth}{!}{\begin{tabular}{p{4.0cm}p{1.9cm}p{1.9cm}p{1.3cm}p{1.6cm}p{1.9cm}p{1.2cm}}
\toprule
Work & Kernel/GP/poly.\ vs neural compared & Tuning or baselines analysed & Ties reported & Selection guidance & Basis check, runs forecast & Open tool \\
\midrule
Hesthaven \& Ubbiali~\cite{hesthaven2018} & -- & -- & -- & -- & -- & -- \\
Guo \& Hesthaven~\cite{guo2018} & -- & -- & -- & -- & -- & -- \\
B\={e}rzi\c{n}\v{s} et al.~\cite{berzins2020} & \checkmark & -- & -- & -- & -- & -- \\
Batlle et al.~\cite{batlle2024} & \checkmark & -- & -- & -- & -- & -- \\
Westermann et al.~\cite{westermann2026} & \checkmark & -- & -- & -- & -- & -- \\
McGreivy \& Hakim~\cite{mcgreivy2024} & -- & \checkmark & -- & -- & -- & -- \\
Williams \& Cremaschi~\cite{williams2019} & -- & -- & -- & \checkmark & -- & -- \\
This paper & \checkmark & \checkmark & \checkmark & \checkmark & \checkmark & \checkmark \\
\bottomrule
\end{tabular}
}\end{table}

\subsection{Motivations}
(i) Method comparisons are reported on fixed problem sets, so a practitioner cannot tell whether a published ranking applies to their problem. (ii) The outcome of a comparison depends on tuning effort~\cite{mcgreivy2024,dodge2019}, so a tool must make the tuning effort explicit and equal. (iii) Small differences are often ties, so a tool should report ties instead of a single winner. (iv) The accuracy of any POD-based model is bounded by the basis (\cref{prop:floor}), which should be checked before tuning a regressor. (v) The practical question is the number of runs needed, which learning-curve extrapolation can answer only roughly and must be validated, not assumed.
